\documentclass[10pt,a4paper]{amsart}
\usepackage[margin=19mm]{geometry}
\usepackage{amsmath,amssymb,mathtools}
\usepackage[scr=rsfs,cal=euler]{mathalfa}
\usepackage{xfrac}
\usepackage[unicode,hidelinks,bookmarks=false]{hyperref}
\newcommand{\ie}{i.e.\ }
\newcommand{\cf}{cf.\ }
\newcommand{\resp}{respectively\ }
\newcommand{\Subsection}{\subsection}
\providecommand{\nobreakdash}{\nobreak\hbox{-}}
\setlength{\emergencystretch}{2em}
\multlinegap0pt
\usepackage{bm}
\usepackage{bbm}
\usepackage{dsfont}
\usepackage{xspace}
\usepackage{cancel}
\usepackage{mathtools}
\usepackage{amsthm}
\makeatletter
\@ifundefined{lemma}{\newtheorem{lemma}{Lemma}}{}
\makeatother
\newcommand{\rset}{\mathbb{R}}
\newcommand{\sphere}{\mathbb{S}}
\title[Polar Depth for Potentially Heavy-Tailed Data]{Polar Depth for Potentially Heavy-Tailed Data}
\author{Stephan Clemen\c{c}on, Carlos Fern\'andes, Pavlo Mozharovskyi, Anne Sabourin}
\date{Source: arXiv:2606.00343v1 (CC BY 4.0)}
\begin{document}
\maketitle

\noindent\emph{Source and licence.} Polar Depth for Potentially Heavy-Tailed Data. Version arXiv:2606.00343v1 (CC BY 4.0), distributed under the Creative Commons Attribution 4.0 International licence, as recorded in the arXiv licence metadata for this exact version and shown on the abstract page https://arxiv.org/abs/2606.00343v1. Full rights notice: licenses/arxiv-2606.00343-CC-BY-4.0.txt.\par\smallskip

\noindent\textit{Editorial scope.} Counted unit: the Lemma whose source label is \texttt{lem:surface\_sphere\_tau\_plusminus\_eps} in the article (source0.tex lines 4134--4145, source label \texttt{lem:surface\_sphere\_tau\_plusminus\_eps}), delivered together with the article's own complete original proof of it (source0.tex lines 4146--4276, one \texttt{proof} environment of 131 lines). No mathematics is written, repaired or re-derived: statement and proof are the article's own text line for line. Required context retained uncounted: eq:surface\_measure\_sphere (lines 1809--1811). Every definition, display and label of those lines is retained unchanged. Omitted from this package and not redistributed: 1--1808, 1812--4133, 4277--4332. Nothing inside a retained excerpt is omitted or altered: every retained body is delivered exactly as the article's lines are written, so no omission receipt of any kind accompanies this package. Citations: the retained excerpts carry no citation command at all, so no bibliography entry of the article is transcribed and no reference is redistributed. Reference adaptation: none was needed and none was made; every reference inside the delivered unit resolves inside it, so no unresolved reference or citation is produced. Printed slips kept literal: no printed word, symbol, hypothesis or number of the counted statement or of its proof was altered. Layout and class: the article's own class is \texttt{article} with packages including authblk, amsthm, amsmath, amsfonts, amssymb, mathtools, mathrsfs, enumitem, hyperref, natbib, color, xcolor, geometry, algorithm; the wrapper is the lane's pinned \texttt{amsart} editorial wrapper at 10pt on A4, which restates the theorem-like environments the retained text needs and adds the article's own macro definitions that the retained text needs (\texttt{\textbackslash rset}, \texttt{\textbackslash sphere}), with extra wrapper packages (bm, bbm, dsfont, xspace, cancel, mathtools). The delivered numbering is the unit's own. Assets: no retained span contains a figure environment, a graphics-inclusion command, a PDF-page inclusion, a table or a TikZ picture; no image file is copied into this package, the package declares no asset dependency and no image file is redistributed with it. Licence: version arXiv:2606.00343v1 of this article is distributed under the Creative Commons Attribution 4.0 International licence, as recorded in the arXiv licence metadata for this exact version and shown on the abstract page https://arxiv.org/abs/2606.00343v1. A full rights notice is delivered at licenses/arxiv-2606.00343-CC-BY-4.0.txt. Characters: every retained excerpt is pure ASCII, so the delivered engine, XeLaTeX with its default Latin Modern text font, reports no missing character.
\begin{equation}\label{eq:surface_measure_sphere}
\sigma_{d-1}(\sphere) = \frac{2 \pi^{d/2} }{\Gamma(d/2)} :=s_{d-1}. 
\end{equation}

    \begin{lemma}[Surface area of $\sphere^{\tau-\varepsilon}\setminus\sphere^{\tau+\varepsilon }$]\label{lem:surface_sphere_tau_plusminus_eps}
      Let $\tau<1/2$ and $\varepsilon \le 1$. Then
$$
      \sigma_{d-1}(\sphere^{\tau-\varepsilon }\setminus\sphere^{\tau+\varepsilon }) \le
      \frac{d^{3/2}}{\sqrt{2\pi}}
      \max\Big(\frac{2 }{\sqrt{1 - 4\tau^2}}, 
       \pi%\frac{\pi}{2\tau}
       \Big)
      \sigma_{d-1}(\sphere_+) \varepsilon.
      $$
    
    \end{lemma}

    \begin{proof}
      Let $\tau<1/2,\epsilon\in(0,1), j\in\{1,\ldots, d\}$ be fixed, and let 
      $M_j = \{x \in \sphere_+: x_j\in
      [\tau-\varepsilon ,\tau+\varepsilon ]\}$. Then
      $$
      M_j = \{x \in \sphere_+: x_j\in
      [\max(\tau-\varepsilon,0) ,\min(\tau+\varepsilon,1) ]\}, 
      $$
      and 
      $\sphere^{\tau-\varepsilon }\setminus\sphere^{\tau+\varepsilon }\subset
      \cup_{j\le d} M_j$,  so that
      %$s^\tau:=
      \begin{equation}
          \label{eq:bound_union_slices1}
          \sigma_{d-1}(\sphere^{\tau-\varepsilon }\setminus\sphere^{\tau+\varepsilon })
      \le d\sigma_{d-1}(M_j).
      \end{equation} 
      Now for fixed $j$, %with $u$ chosen as
      %the $j^{th}$ canonical basis vector,  % and the notations from the proof of Lemma~\ref{lem:surface_spherical_caps}, %$M_j$ is the
      let $\Delta_j$ denote the subset of $\sphere$ comprised between the parallel circles
      $\mathcal{C}_{\max(\tau-\varepsilon,0) }$ and $\mathcal{C}_{\min(\tau+\varepsilon, 1) }$, where $\mathcal{C}_b = \{x\in\sphere: x_j = b\}$. Then $M_j = \Delta_j\cap\sphere_+$, and by symmetry with respect to the $d-1$ coordinates different from $j$, we have 
      \begin{equation}
          \label{eq:bound_union_slices_2}
          \sigma_{d-1}(M_j) = 2^{-d+1}\sigma_{d-1}(\Delta_j).
      \end{equation} 
      We are left with computing the surface measure $\sigma_{d-1}(\Delta_j)$. 
      For $0<b<1$, the circle $\mathcal{C}_b$ is isometric with the sphere in the ambient space $\rset^{d-1}$ with radius $r_b = \sqrt{1-b^2}\le 1$. % It is known
Recall from~\eqref{eq:surface_measure_sphere} that       that the surface of the $d-2$-dimensional unit sphere is 
       $s_{d-2} = \frac{2 \pi^{(d-1)/2}}{\Gamma((d-1)/2)},  $
so that % , with $s_{d-1} = \sigma_{d-1}(\sphere)$, we have 
      \begin{align*}
        \frac{s_{d-2}}{s_{d-1}}  = \frac{\pi^{-1/2} \Gamma(d/2)}{\Gamma((d-1)/2)}
        \le \sqrt{\frac{d}{2\pi}}, 
      \end{align*}
      where we have used Gautschi's inequality to bound from above the ratio of Gamma's. 
      Whence, for $0\le b\le 1$, we have
$$
\sigma_{d-2}(\mathcal{C}_b) = r_b^{d-2} s_{d-2} = 
\le \sqrt{\frac{d}{2\pi}} \sigma_{d-1}(\sphere).%,   % \frac{2 \pi^{(d-1)/2}}{\Gamma((d-1)/2)}
$$
     
      % By symmetry of the problem with
      % respect to the $d-1$ coordinates different from $j$, intersecting
      % with $\sphere_+$ amounts to dividing the surface by
      % $2^{d-1}$. 
      Finally, for $\varepsilon<\tau <1/2$, the geodesic distance between the two circles $\mathcal{C}_{\tau- \varepsilon }$ and $\mathcal{C}_{\tau+\varepsilon }$ is
      $$
      \rho(\mathcal{C}_{\tau-\varepsilon }, \mathcal{C}_{\tau+\varepsilon }) =
      \arcsin(\tau+\varepsilon ) - \arcsin(\tau-\varepsilon ) \le
      \frac{2 \varepsilon }{\sqrt{1 - (\tau+\varepsilon )^2}} \le 
      \frac{2 \varepsilon }{\sqrt{1 - 4\tau^2}}.
      $$
      Also for $\varepsilon\ge \tau$, 
$$
      \rho(\mathcal{C}_{0}, \mathcal{C}_{\min(\tau+\varepsilon,1) }) \le 
      \rho(\mathcal{C}_{0}, \mathcal{C}_{\min(2\varepsilon,1) }) \le 
      \arcsin(\min(2\varepsilon,1)) \le \pi\varepsilon,
      %\le \frac{\pi/2}{\tau} \varepsilon.
      $$
      where the last inequality is obtained by studying the variations of the functions $x\mapsto \arcsin(\min(2x,1))/x$ on $(0,1)$. 
      Combining the latter  three displays, we obtain 
      \begin{equation}\label{eq:bound_union_slice_3}
      \sigma_{d-1}(\Delta_j) \le \sup_{0\le b\le 1}\sigma_{d-1}(\mathcal{C}_b)\rho(\mathcal{C}_{\max(\tau-\varepsilon,0) }, \mathcal{C}_{\min(\tau+\varepsilon,1) }) 
      \le \sqrt{\frac{d}{2\pi}} \sigma_{d-1}(\sphere)
      \max\Big(\frac{2 }{\sqrt{1 - 4\tau^2}}, 
       {\pi}\Big)\varepsilon.
      %\frac{2 \varepsilon }{\sqrt{1 - (\tau+\varepsilon )^2}}, 
      \end{equation}
      Finally,  combining Equations~\eqref{eq:bound_union_slices1}, \eqref{eq:bound_union_slices_2}, \eqref{eq:bound_union_slice_3} and using $\sigma(\sphere_+) = 2^{-d}\sigma_{d-1}(\sphere)$, we obtain
      \begin{align*}
      \sigma_{d-1}(\sphere^{\tau-\varepsilon}\setminus\sphere^{\tau+\varepsilon}) 
      & \le 
      d \sigma_{d-1}(M_j) \le d \frac{\sigma_{d-1}(\Delta_j)}{2^{d-1}}\\
      &\le \frac{d^{3/2}}{\sqrt{2\pi}}
      \max\Big(\frac{2 }{\sqrt{1 - 4\tau^2}}, 
       \pi %\frac{\pi}{2\tau}
       \Big)
      \frac{\sigma_{d-1}(\sphere)}{2^{d-1}} \varepsilon\\
      &=\frac{d^{3/2}}{\sqrt{2\pi}}
      \max\Big(\frac{2 }{\sqrt{1 - 4\tau^2}}, 
       \pi%\frac{\pi}{2\tau}
       \Big)
      \sigma_{d-1}(\sphere_+) \varepsilon.\\
      %&\le \frac{(2d)^{3/2}}{\sqrt{\pi}}
      %\frac{\varepsilon }{\sqrt{1 - (\tau+\varepsilon )^2}}
      %\sigma_{d-1}(\sphere_+).
      \end{align*}
      % the greatest circle $\mathcal{C}_{\tau-\varepsilon }$
      % has radius $r=\sqrt{1 - (\tau-\varepsilon )^2}\le 1 $.  Thus, by the
      % same integral
      %With the same argument as in the proof of
      %Lemma~\ref{lem:surface_spherical_caps}, with $s_{d-2}$ denoting
      %the surface area of the $d-2$ sphere in $\rset^{d-1}$, we obtain 
      % $$
      % \sigma_{d-1}(M_j) ~\le~
      % \frac{1}{2^{d-1}} \frac{2 \varepsilon }{\sqrt{1 - (\tau+\varepsilon )^2}}  s_{d-2}
      % % \int_{b=\tau-\varepsilon }^{\tau+\varepsilon }  s_{d-2} \ud b
      % % ~ = ~  \frac{2\varepsilon }{2^{d-1}} s_{d-2}
      % % ~
      % % \le ~\frac{ \varepsilon }{2^{d-2}} \sqrt{\frac{d}{2\pi}} \sigma_{d-1}(\sphere),
      % ~\le ~
      % \frac{1}{2^{d-1}} \frac{2 \varepsilon }{\sqrt{1 - (\tau+\varepsilon )^2}} 
      % \sqrt{\frac{d}{2}} \sigma_{d-1}(\sphere), 
      % $$
%where the latter inequality derives from Gautschi's inequality. 

%       Consider the great circle
%       $\mathcal{C}_0 = \{x \in \sphere:  \langle x, u \rangle = 0 \}$ and the parallel circle at height $\beta$, 
%       $\mathcal{C}_\beta = \{x \in \sphere:  \langle x, u \rangle = \beta \}$.
%       Let $\rho_\beta>0$ denote the geodesic distance between the two. 
%Let $s_{d-2}$ denote the $d-2$ dimensional surface area of the great circle $\mathcal{C}_0$, which is isometric to the $d-2$ sphere. 
%Then the surface measure of $M$ is less than the product $\rho_\beta \times s_{d-2}$. On the one hand,  $\rho_\beta = \arcsin\beta \le \frac{\beta}{\sqrt{1-\beta^2}}$.
%       %       Then for $0\le b \le \beta$,
% %       $\mathcal{C}_b$ is isometric to the $d-2$-sphere in the ambient space $\rset^{d-1}$, with radius $r(b) = \sqrt{1- b^2}$. Let  $s_{d-2}$ denote the surface area of $\mathcal{C}_0$.  With these notations the surface area of $\mathcal{C}_b$ is $s_{d-1}r(b)^{d-2}$ and 
% %        $$\sigma_{d-1}(M)  = \int_{b=0}^\beta  s_{d-2}r(b)^{d-1} \ud b \le  s_{d-2} \beta. $$
% %    %    On the one hand, 
% % %       $$
% % % \rho_\beta = \arcsin(\beta) \le \beta/\sqrt{1- \beta^2}.
% % %       $$
%       On the other hand  it is known that the surface of the $d-2$-dimensional sphere is 
%       $$s_{d-2} = \frac{2 \pi^{(d-1)/2}}{\Gamma((d-1)/2)},  $$
%       
%       so that, with $s_{d-1} = \sigma_{d-1}(\sphere)$, we have 
%       \begin{align*}
%         \frac{s_{d-2}}{s_{d-1}}  = \frac{\pi^{-1/2} \Gamma(d/2)}{\Gamma((d-1)/2)}
%         \le \sqrt{\frac{d}{2\pi}}, 
%       \end{align*}
%       where we have used Gautschi's inequality to bound from above the ratio of Gamma's. 

      
    \end{proof}

\end{document}
